\documentclass[12pt, twoside]{article}
%\documentclass[12pt, twoside]{article}
\usepackage[letterpaper, margin=1in, headsep=0.2in]{geometry}
\setlength{\headheight}{0.6in}
%\usepackage[english]{babel}
\usepackage[utf8]{inputenc}
\usepackage{microtype}
\usepackage{amsmath}
\usepackage{amssymb}
%\usepackage{amsfonts}
\usepackage[nomessages]{fp} %\FPeval{\var-name}{2*sin(pi/6)}
\usepackage{siunitx} %units in math. eg 20\milli\meter
\usepackage{yhmath} % for arcs, overparenth command
\usepackage{tikz} %graphics
\usetikzlibrary{quotes, angles, arrows, arrows.meta}
\usepackage{pgfplots}
\usepgfplotslibrary{fillbetween}
\pgfplotsset{compat=1.16}
\usepackage{graphicx} %consider setting \graphicspath{{images/}}
\usepackage{parskip} %no paragraph indent
\usepackage{enumitem}
\usepackage{multicol}
\usepackage{venndiagram}

\usepackage{fancyhdr}
\pagestyle{fancy}
\fancyhf{}
\renewcommand{\headrulewidth}{0pt} % disable the underline of the header
\raggedbottom
\hfuzz=2mm %suppresses overfull box warnings

\usepackage{hyperref}

\fancyhead[LO]{La Scuola d'Italia / Dr.\ Huson / IB Math 11 \\* 5 October 2026}
\fancyhead[RO]{\fbox{\begin{minipage}{6cm}\footnotesize Nickname:\\[0.5cm]\end{minipage}}}
\fancyhead[LE]{\fbox{\begin{minipage}{6cm}\footnotesize Nickname:\\[0.5cm]\end{minipage}}}
\fancyfoot[C]{\thepage}

\begin{document}
\subsubsection*{1.9 Pre-Quiz: Arithmetic and Geometric Sequences}

\emph{Review for Wednesday's quiz. A calculator is allowed. Show your
working: write the formula you use, then substitute.}

\begin{enumerate}[itemsep=0.15cm]

%--- Problem 1: Arithmetic sequence; d, u_n, u_40, which values are terms ---
%--- Source: Original (freshly authored, AA SL 1.2); review of 1-1CW, 1-3CW ---
\item[1.]

Consider the arithmetic sequence
$$12,\; 19,\; 26,\; 33,\; \ldots$$

\begin{enumerate}[itemsep=0.1cm]
  \item Write down the value of the common difference $d$.
  \item Find an expression for $u_{n}$, the $n^{\text{th}}$ term.
  \item Find $u_{40}$.
  \item Show that $250$ is a term of the sequence, and find which term
  it is.
  \item Determine whether $300$ is a term of the sequence. Justify your
  answer.
\end{enumerate}

\begin{tikzpicture}
  \draw (-0.6,0) rectangle (14.6,14.8);
  \foreach \y in {1,2,...,14} \draw[dotted] (0,\y)--(14.1,\y);
\end{tikzpicture}

\newpage

%--- Problem 2: Two given terms -> u_1 and d; first negative term; terms in k ---
%--- Source: Original (freshly authored, AA SL 1.2); review of 1-2CW, 1-3CW ---
\item[2.]

In an arithmetic sequence, $u_{5} = 31$ and $u_{12} = 3$.

\begin{enumerate}[itemsep=0.1cm]
  \item Find the common difference $d$ and the first term $u_{1}$.
  \item Find an expression for $u_{n}$.
  \item Find the first term of the sequence that is negative.
\end{enumerate}

The first three terms of a different arithmetic sequence are
$$k + 2, \qquad 3k, \qquad 4k + 1.$$

\begin{enumerate}[itemsep=0.1cm]
  \item[(d)] Find the value of $k$, and write down the first three terms.
  \item[(e)] Find the $20^{\text{th}}$ term of this sequence.
\end{enumerate}

\begin{tikzpicture}
  \draw (-0.6,0) rectangle (14.6,14.3);
  \foreach \y in {1,2,...,14} \draw[dotted] (0,\y)--(14.1,\y);
\end{tikzpicture}

\newpage

%--- Problem 3: Arithmetic, geometric or neither; geometric u_n, u_10, first term over 1000; two given terms ---
%--- Source: Original (freshly authored, AA SL 1.2, 1.3); review of 1-6CW ---
\item[3.]

\begin{enumerate}[itemsep=0.1cm]
  \item State whether each sequence is arithmetic, geometric or neither.
  Write down $d$ or $r$ where there is one.

  \hspace*{1em} (i) $3,\; 7,\; 11,\; 15$ \hfill (ii) $3,\; 6,\; 12,\; 24$
  \hfill (iii) $3,\; 4,\; 6,\; 9$ \hfill (iv) $80,\; 20,\; 5,\; 1.25$
  \hspace*{1em}
\end{enumerate}

Consider the geometric sequence
$$6,\; 9,\; 13.5,\; \ldots$$

\begin{enumerate}[itemsep=0.1cm]
  \item[(b)] Find the common ratio $r$.
  \item[(c)] Write down an expression for $u_{n}$.
  \item[(d)] Find $u_{10}$.
  \item[(e)] Find the smallest value of $n$ for which $u_{n} > 1000$.
  \item[(f)] A different geometric sequence has $u_{2} = 12$ and
  $u_{5} = 324$. Find the common ratio and the first term.
\end{enumerate}

\begin{tikzpicture}
  \draw (-0.6,0) rectangle (14.6,12.6);
  \foreach \y in {1,2,...,12} \draw[dotted] (0,\y)--(14.1,\y);
\end{tikzpicture}

\newpage

%--- Problem 4: Applications -- arithmetic vs geometric salary; compound interest and real value ---
%--- Source: Original (freshly authored, AA SL 1.2, 1.3, 1.4); review of 1-7CW, 1-8CW ---
\item[4.]

Elena is offered two jobs.

\begin{center}
\begin{tabular}{rl}
Job A: & \$$42\,000$ in the first year, rising by \$$1500$ each year. \\[2pt]
Job B: & \$$40\,000$ in the first year, rising by $5\%$ each year. \\
\end{tabular}
\end{center}

\begin{enumerate}[itemsep=0.1cm]
  \item Find Elena's salary in year $3$ for each job.
  \item Write down an expression for the salary in year $n$ for each job.
  \item Find the first year in which the Job B salary is greater than the
  Job A salary.
\end{enumerate}

Elena invests \$$3000$ at $3\%$ per year compounded annually.

\begin{enumerate}[itemsep=0.1cm]
  \item[(d)] Find the value of the investment after $8$ years.
  \item[(e)] Inflation is $2\%$ per year. Find the real value of the
  investment after $8$ years.
\end{enumerate}

\begin{tikzpicture}
  \draw (-0.6,0) rectangle (14.6,12.9);
  \foreach \y in {1,2,...,12} \draw[dotted] (0,\y)--(14.1,\y);
\end{tikzpicture}

\end{enumerate}
\end{document}
